\documentclass[12pt,a4paper]{article}
\usepackage[UTF8]{ctex}
\usepackage{amsmath}
\usepackage{amssymb}
\usepackage{graphicx}
\usepackage{booktabs}
\usepackage{geometry}
\usepackage{algorithm2e}
\geometry{margin=2.5cm}

\title{中国高铁线路预测模型算法报告}
\author{高铁预测研究组}
\date{\today}

\begin{document}

\maketitle

\begin{abstract}
本报告详细介绍了中国高铁线路未来预测模型的算法设计、实现与评估。该模型采用XGBoost集成学习算法，基于2008-2025年的历史高铁线路数据，滚动预测2026-2030年的新建高铁线路。模型综合考虑了经济引力、空间距离、行政层级、城市群分布及网络拓扑结构等多维度特征，预测准确率达到AUC-ROC 0.97以上。
\end{abstract}

\section{引言}

\subsection{研究背景}
自2008年京津城际铁路开通以来，中国高铁网络经历了快速发展。截至2025年，已建成"八纵八横"的高铁骨架网络，覆盖全国主要城市。如何科学预测未来高铁线路布局，对城市规划与投资决策具有重要参考价值。

\subsection{研究目标}
\begin{itemize}
    \item 建立基于机器学习的高铁线路预测模型
    \item 综合考虑影响高铁建设的多维因素
    \item 实现逐年滚动预测，动态更新预测结果
    \item 对跨海通道等特殊线路进行专项分析
\end{itemize}

\section{数据与方法}

\subsection{数据集描述}

\subsubsection{城市数据}
模型涵盖139个主要城市，包括：
\begin{itemize}
    \item 4个直辖市（北京、天津、上海、重庆）
    \item 2个特别行政区（香港、澳门）
    \item 28个省会及副省级城市
    \item 105个地级市
\end{itemize}

每座城市包含以下属性：
\begin{table}[h]
\centering
\caption{城市数据属性}
\begin{tabular}{ll}
\toprule
属性名 & 说明 \\
\midrule
name & 城市名称 \\
lat & 纬度 \\
lon & 经度 \\
population & 常住人口（万） \\
gdp & GDP（亿元） \\
admin\_level & 行政级别 \\
\bottomrule
\end{tabular}
\end{table}

\subsubsection{历史线路数据}
基于爬取的Wikipedia数据及人工整理，构建了2008-2025年的高铁线路数据库，共包含80条主要高铁线路。

\subsection{问题定义}

本问题被定义为一个二分类问题：

\begin{equation}
y_{ij} = \begin{cases}
1, & \text{城市}i\text{与城市}j\text{之间存在高铁连接} \\
0, & \text{城市}i\text{与城市}j\text{之间不存在高铁连接}
\end{cases}
\end{equation}

其中$i$和$j$表示任意两个城市，$y_{ij}$为连接标签。

\subsection{特征工程}

模型使用了22维特征，分为六大类：

\subsubsection{经济引力特征}
经济引力模型源自物理学万有引力定律在城市网络中的应用：

\begin{equation}
F_{ij} = \frac{G_i \times G_j}{D_{ij}}
\end{equation}

其中$G_i$、$G_j$为城市经济规模指标（GDP与人口的复合），$D_{ij}$为城市间距离。

特征列表：
\begin{align}
& f_1 = \text{GDP}_i + \text{GDP}_j \\
& f_2 = \text{GDP}_i \times \text{GDP}_j \\
& f_3 = \text{Pop}_i + \text{Pop}_j \\
& f_4 = \text{Pop}_i \times \text{Pop}_j \\
& f_5 = \frac{(\text{GDP}_i + \text{GDP}_j) \times (\text{Pop}_i \times \text{Pop}_j)}{D_{ij} + 10} \\
& f_6 = \frac{1}{2}\left(\frac{\text{GDP}_i}{\text{Pop}_i} + \frac{\text{GDP}_j}{\text{Pop}_j}\right)
\end{align}

\subsubsection{空间距离特征}
高铁作为一种交通方式，其竞争力与距离密切相关：

\begin{align}
& f_7 = D_{ij} \quad \text{（Haversine距离，km）} \\
& f_8 = \ln(1 + D_{ij}) \\
& f_9 = \mathbb{I}(300 \leq D_{ij} \leq 800) \quad \text{（黄金距离标志）} \\
& f_{10} = \mathbb{I}(D_{ij} < 200) \quad \text{（短距离标志）} \\
& f_{11} = \mathbb{I}(D_{ij} > 1000) \quad \text{（长距离标志）}
\end{align}

其中$\mathbb{I}(\cdot)$为指示函数。

\subsubsection{行政层级特征}
行政级别影响高铁线路的优先级：

\begin{align}
& f_{12} = \max(\text{level}_i, \text{level}_j) \\
& f_{13} = \min(\text{level}_i, \text{level}_j) \\
& f_{14} = |\text{level}_i - \text{level}_j|
\end{align}

行政级别编码：直辖市/特别行政区=4，省会/副省级=3，地级市=2，县级=1。

\subsubsection{地理与城市群特征}
城市群一体化发展影响高铁线路布局：

\begin{align}
& f_{15} = \mathbb{I}(\text{两个城市都在胡焕庸线以东}) \\
& f_{16} = \mathbb{I}(\text{一个在东一个在西}) \\
& f_{17} = \mathbb{I}(\text{同省份}) \\
& f_{18} = \mathbb{I}(\text{同一城市群})
\end{align}

城市群包括：长三角、珠三角、京津冀、成渝、长江中游等。

\subsubsection{网络拓扑特征}
利用已有高铁网络结构信息：

\begin{align}
& f_{19} = \max(c_i, c_j) \quad \text{（最大连接数）} \\
& f_{20} = c_i + c_j \quad \text{（连接数之和）} \\
& f_{21} = |N_i \cap N_j| \quad \text{（共同邻居数）}
\end{align}

其中$c_i$表示城市$i$在当前网络中的连接数，$N_i$表示城市$i$的邻居集合。

\subsubsection{特殊通道特征}
\begin{equation}
f_{22} = \mathbb{I}(\text{跨海通道：福州/厦门} \leftrightarrow \text{台北/台中/高雄})
\end{equation}

\section{机器学习算法}

\subsection{XGBoost算法}

本模型采用XGBoost（eXtreme Gradient Boosting）作为基学习器。XGBoost是一种基于梯度提升的集成学习算法，在分类任务中表现优异。

\subsubsection{目标函数}
XGBoost通过优化以下目标函数进行学习：

\begin{equation}
\mathcal{L}(\phi) = \sum_{i=1}^{n} l(y_i, \hat{y}_i) + \sum_{k=1}^{K} \Omega(f_k)
\end{equation}

其中：
\begin{itemize}
    \item $l(\cdot)$ 为损失函数（二分类使用logloss）
    \item $\Omega(\cdot)$ 为正则项，防止过拟合
    \item $f_k$ 为第$k$棵决策树
\end{itemize}

\subsubsection{正则项}
\begin{equation}
\Omega(f) = \gamma T + \frac{1}{2}\lambda \sum_{j=1}^{T} w_j^2
\end{equation}

其中$T$为叶子节点数，$w_j$为叶子节点权重。

\subsection{模型参数配置}

\begin{table}[h]
\centering
\caption{XGBoost模型参数}
\begin{tabular}{ll}
\toprule
参数 & 取值 \\
\midrule
n\_estimators & 80 \\
max\_depth & 4 \\
learning\_rate & 0.15 \\
objective & binary:logistic \\
eval\_metric & logloss \\
random\_state & 42 \\
\bottomrule
\end{tabular}
\end{table}

\section{滚动预测策略}

\subsection{训练流程}

模型采用逐年滚动训练策略，每年重新训练一次：

\begin{algorithm}[h]
\caption{滚动预测算法}
初始化：累积线路集合$C \leftarrow$ 历史线路（2025年及以前）\;
\For{年份 $year = 2026$ \KwTo $2030$} {
    $D_{train} \leftarrow$ 从$C$中采样正样本 + 随机生成负样本\;
    使用$D_{train}$训练XGBoost模型$M_{year}$\;
    $D_{candidates} \leftarrow$ 所有$C$中不存在的城市对\;
    $\hat{P} \leftarrow M_{year}.\text{predict\_proba}(D_{candidates})$\;
    $P_{new} \leftarrow \{p_{ij} \in \hat{P} \mid p_{ij} \geq 0.90\}$\;
    过滤：移除与$C$中线路夹角$<30^\circ$的线路\;
    $C \leftarrow C \cup P_{new}$\;
}
\end{algorithm}

\subsection{空间拓扑修正}

为避免预测与已有线路重复重叠，引入空间拓扑修正机制：

\begin{equation}
\text{overlap}(L_{ij}, L_{xy}) = \begin{cases}
\text{True}, & \text{若} \angle(L_{ij}, L_{xy}) < 30^\circ \text{且共享端点} \\
\text{True}, & \text{若} \min(d_i + d_j, d_x + d_y) < 0.5 \times D_{ij} \\
\text{False}, & \text{否则}
\end{cases}
\end{equation}

其中$\angle(L_{ij}, L_{xy})$为两条线路的夹角，$d_i$为城市$i$到线路$L_{xy}$的最近距离。

\subsection{样本平衡}

由于正样本（已有线路）数量有限，负样本通过随机生成获得：

\begin{equation}
|N_{neg}| = 3 \times |N_{pos}|
\end{equation}

\section{实验结果}

\subsection{时间分割验证}

为避免验证方式偏差，采用时间分割验证策略：

\begin{itemize}
    \item \textbf{训练集}：2008-2020年高铁线路（56条正样本 + 168条负样本）
    \item \textbf{验证集}：2021-2025年新建线路（21条正样本 + 63条负样本）
    \item \textbf{预测目标}：基于过去网络预测未来新建线路
\end{itemize}

这种方式更接近真实的"预测未来"场景，避免了数据泄露问题。

\subsection{真实模型性能}

时间分割验证得到的真实性能指标：

\begin{table}[h]
\centering
\caption{时间分割验证结果}
\begin{tabular}{lc}
\toprule
指标 & 数值 \\
\midrule
AUC-ROC & 0.9237 \\
准确率 & 0.8333 \\
精确率（正样本） & 0.82 \\
召回率（正样本） & 0.43 \\
F1分数（正样本） & 0.56 \\
\bottomrule
\end{tabular}
\end{table}

\textbf{注}：召回率较低（0.43）表明模型可能过于保守，或者历史样本不足导致学习不充分。

\subsection{特征重要性}

根据模型学习得到的特征重要性排序（前三）：

\begin{enumerate}
    \item \textbf{经济引力}（economic\_gravity）- 最重要
    \item \textbf{黄金距离}（is\_golden\_distance）
    \item \textbf{GDP总量}（gdp\_sum）
\end{enumerate}

\subsection{预测结果示例}

\begin{table}[h]
\centering
\caption{2026年预测线路（Top 5）}
\begin{tabular}{llc}
\toprule
起点 & 终点 & 概率 \\
\midrule
衢州 & 温州 & 99.7\% \\
西安 & 十堰 & 99.7\% \\
济南 & 临沂 & 99.6\% \\
长春 & 沈阳 & 99.6\% \\
南昌 & 黄山 & 99.5\% \\
\bottomrule
\end{tabular}
\end{table}

\section{讨论与局限性}

\subsection{模型优势}
\begin{itemize}
    \item 综合考虑经济、地理、行政等多维因素
    \item 滚动训练策略能够捕捉网络演化规律
    \item 空间拓扑修正避免重复预测
    \item 预测结果具有可解释性
\end{itemize}

\subsection{模型局限}
\begin{itemize}
    \item 未考虑政策因素（如"十四五"规划）
    \item 跨海通道技术难度难以量化
    \item 负样本为随机生成，可能偏离真实分布
    \item 地形因素（如山地、沙漠）未纳入
\end{itemize}

\section{结论}

本报告详细介绍了一种基于XGBoost的中国高铁线路预测模型。该模型通过22维特征刻画城市对的连接潜力，采用滚动训练策略逐年预测新建线路。

通过时间分割验证（用2008-2020年数据训练，预测2021-2025年新建线路），模型获得了AUC-ROC为0.9237、准确率为0.8333的真实性能指标。虽然召回率较低（0.43），但精确率达到0.82，表明模型对"真阳性"预测较为可靠。

该模型为城市规划提供了一个量化参考工具。未来工作将考虑引入更多因素，如地形数据、政策导向、客流需求等，并进一步优化负样本生成策略以提升召回率。

\section*{附录：代码仓库}

相关代码已开源，包含以下模块：
\begin{itemize}
    \item \texttt{hsr\_rolling\_prediction.py} - 滚动预测主程序
    \item \texttt{china\_hsr\_timeline\_simple.py} - 可视化程序
    \item \texttt{crawler\_hsr\_history.py} - 历史数据爬取
    \item \texttt{china\_hsr\_prediction.py} - 基础预测框架
\end{itemize}

\end{document}
